% Part: normal-modal-logic
% Chapter: tableaux
% Section: simple-S5

\documentclass[../../../include/open-logic-section]{subfiles}

\begin{document}

\olfileid[tr]{nml}{tab}{s5}

\olsection{\Log{S5} için Basit \usetoken{P}{tableau}}

\Log{S5}, evrensel modeller sınıfına, yani her dünyadan her dünyaya
erişilebilen modeller sınıfına göre sağlam ve tamdır. Evrensel modellerde
erişilebilirlik bağıntısının ayrıntısı önem taşımaz: ``$\mSat{M}{!A}[w]$
olan bir $w$ dünyası vardır'' ifadesi, ancak ve ancak $u$'dan erişilebilen
böyle bir $w$ varsa doğrudur. Dolayısıyla \Log{S5} içinde modelleri yalnızca
bir dünyalar kümesi ve bir $V$ değerlemesi olarak tanımlayabiliriz. Bu,
!!{tableau} kurallarını da sadeleştirebilmemiz gerektiğini düşündürür. Genel
durumda önek olarak pozitif tam sayı dizileri alırız; böylece hangi öneklerin
başkalarından erişilebilen dünyaları adlandırdığını izleyebiliriz: $\sigma.n$,
$\sigma$'dan erişilebilen bir dünyayı adlandırır. Fakat \Log{S5} içinde her
dünyadan her dünyaya erişilebildiğinden bunu izlemeye gerek yoktur. Bunun
yerine pozitif tam sayıları önek olarak kullanabiliriz. Sadeleştirilmiş
kurallar \olref{tab:rules-S5}'te verilmiştir.

\begin{table}
  \begin{center}
    \def\arraystretch{3}\def\fCenter{}
    \begin{tabular}{|c|c|}
    \hline
    \iftag{prvBox}{
      \AxiomC{\sFmla{\True}{\Box !A}[n]}
      \RightLabel{\TRule{\True}{\Box}}
      \UnaryInfC{\sFmla{\True}{!A}[m]}
      \DisplayProof
      &
      \AxiomC{\sFmla{\False}{\Box !A}[n]}
      \RightLabel{\TRule{\False}{\Box}}
      \UnaryInfC{\sFmla{\False}{!A}[m]}
      \DisplayProof\\
      $m$ kullanılmıştır & $m$ yenidir
      \\[1ex]
      \hline}{}
    \iftag{prvDiamond}{
      \AxiomC{\sFmla{\True}{\Diamond !A}[n]}
      \RightLabel{\TRule{\True}{\Diamond}}
      \UnaryInfC{\sFmla{\True}{!A}[m]}
      \DisplayProof
      &
      \AxiomC{\sFmla{\False}{\Diamond !A}[n]}
      \RightLabel{\TRule{\False}{\Diamond}}
      \UnaryInfC{\sFmla{\False}{!A}[m]}
      \DisplayProof\\
      $m$ yenidir & $m$ kullanılmıştır
      \\[1ex]
      \hline}{}
    \end{tabular}
  \end{center}
  \caption{\Log{S5} için sadeleştirilmiş kurallar.}
  \ollabel{tab:rules-S5}
\end{table}

\begin{ex}
  $\Log{S5} \Proves \Ax{5}$, yani $\Diamond!A \lif \Box\Diamond!A$
  olduğunu gösteren sadeleştirilmiş kapalı bir tablo veriyoruz.
  \begin{oltableau}
    [\pFmla{\False}{\Diamond\formula{A} \lif \Box\Diamond \formula{A}}{1},
      just = \TAss
      [\pFmla{\True}{\Diamond \formula{A}}{1}, just = {\TRule{\False}{\lif}[1]}
        [\pFmla{\False}{\Box\Diamond \formula{A}}{1},
          just = {\TRule{\False}{\lif}[1]}
          [\pFmla{\False}{\Diamond \formula{A}}{2},
            just = {\TRule{\False}{\Box}[3]}
            [\pFmla{\True}{\formula{A}}{3},
              just = {\TRule{\True}{\Diamond}[2]}
                [\pFmla{\False}{\formula{A}}{3},
                  just = {\TRule{\False}{\Diamond}[4]}, close]
            ]
          ]
        ]
      ]
    ]
  \end{oltableau}
\end{ex}

\end{document}
